\chapter{电力电子技术的应用}

前九章讲的是器件与电路，本章讲\textbf{装置}：把器件、电路、控制组合起来，解决实际的工程问题。教材按应用领域分了七节，主线是“电动机传动（最大市场）→ 电源（UPS、开关电源、PFC）→ 电力系统 → 其他”。

% ============================================================
\section{晶闸管直流电动机系统}

晶闸管可控整流装置带直流电动机负载组成的系统，习惯称为\textbf{晶闸管直流电动机系统}，是直流调速的主要形式。直流电动机的电枢回路有电阻 $R$、电感 $L$ 和反电动势 $E$，$E=C_e\Phi n$。

\subsection{工作于整流状态时}

电流连续时，整流输出平均电压 $U_d$ 与电动机反电动势平衡：

\keyformula{U_d=E+I_dR_\Sigma=C_e\Phi n+I_dR_\Sigma}

其中 $R_\Sigma$ 为电枢电阻与整流装置等效内阻（含换相压降折算）之和。把 $U_d=2.34U_2\cos\alpha-\Delta U_d$ 代入，得\textbf{机械特性}：

\keyformula{n=\frac{2.34U_2\cos\alpha-\Delta U_d-I_dR_\Sigma}{C_e\Phi}}

\begin{itemize}
  \item 改变 $\alpha$ 即改变理想空载转速 $n_0$，得到一族\textbf{平行下斜的直线}——这就是晶闸管直流调速的调压调速原理；
  \item 与直流发电机供电相比，晶闸管系统的机械特性斜率更大（等效内阻大），且\textbf{电流断续时特性明显变软}（$n$ 随 $I_d$ 减小而迅速升高，理想空载转速虚高），这是它特有的现象，分析时必须先判断电流是否连续。
\end{itemize}

\subsection{工作于有源逆变状态时}

$\alpha>90^\circ$（即逆变角 $\beta<90^\circ$）时 $U_d<0$，电动机作\textbf{发电机运行}（回馈制动），能量经变流器送回电网。此时

\keyformula{n=\frac{-2.34U_2\cos\beta-\Delta U_d-I_dR_\Sigma}{C_e\Phi}}

负号表示\textbf{逆变时电动机的转向与整流时相反}（相对于同一 $I_d$ 方向）。对应不同的 $\beta$ 可得到一族平行的机械特性，位于第二象限（正转制动）或第四象限（反转制动）。

\subsection{直流可逆电力拖动系统}

要实行四象限运行（正转电动、正转制动、反转电动、反转制动），须使电动机的\textbf{电压和电流都能反向}，常用两套变流装置\textbf{反并联}供电：

\begin{itemize}
  \item \textbf{有环流接线}：两组同时加触发脉冲，用\textbf{环流电抗器}限制两组之间的环流；切换平滑、无死区，但需笨重电抗器、损耗大（如三相半波有环流接线）。
  \item \textbf{无环流接线}：任何时刻只让一组工作，另一组封锁脉冲；切换时须先封锁、等电流过零并留死区，再开放另一组，故\textbf{切换有死区}，但效率高、无需环流电抗器（如三相全控桥无环流接线）——实际多用。
\end{itemize}

% ============================================================
\section{变频器和交流调速系统}

过去调速传动的主流是晶闸管直流电动机系统，但直流电动机有电刷与换向器、维护麻烦、容量受限。随着变频技术发展，\textbf{交流电动机的调速性能已可与直流电动机相媲美}，交流调速逐渐占据主导地位。

\subsection{交—直—交变频器}

\textbf{交—直—交变频器}（间接交流变流电路）由“整流器 + 中间直流环节 + 逆变器”构成，是通用变频器的主流结构。按中间直流环节的储能元件分为电压型与电流型。

\begin{itemize}
  \item \textbf{电压型间接交流变流电路}：中间为大电容，逆变器为电压型 PWM 逆变。若整流部分用不可控二极管整流，则\textbf{不能处理再生反馈的电力}（能量只能单方向流动），制动时需加制动电阻或制动单元把能量消耗掉，否则直流母线电压会升高（“泵升电压”）。
  \item \textbf{可逆的电压型电路}：整流和逆变均为 PWM 控制（双 PWM 变频器），可实现能量的双向流动，能再生制动、且输入电流正弦、功率因数可调。
  \item \textbf{电流型间接交流变流电路}：中间为大电感，可方便地实现再生反馈（通过改变整流桥的极性或采用可逆整流），但动态响应较慢。
\end{itemize}

\subsection{交—交变频器}

\textbf{交—交变频器}（直接变频、周波变流器）不经中间直流环节，用两组反并联相控整流直接把电网频率变成较低频率（见第 6 章）。其特点是\textbf{效率高、低频输出波形好、可四象限运行}，适合\textbf{低速大功率}传动（轧机、矿井提升机、水泥磨机）；缺点是输出频率不能超过输入频率的 $1/3\sim1/2$、功率因数低、输入谐波大。

\subsection{变频调速的基本控制方式}

由 $n_0=\dfrac{60f_1}{p}$ 可知，改变频率 $f_1$ 即可改变同步转速。但只变频不变压会引起磁通变化（$U_1\approx E_1=4.44f_1N_1\Phi_m$），因此须\textbf{变压变频（VVVF）}配合：

\keyformula{\frac{U_1}{f_1}=\text{const}\quad\Longrightarrow\quad \Phi_m\approx\text{const}}

\begin{itemize}
  \item \textbf{基频以下}：保持 $U_1/f_1=$ 常数（恒磁通），实现\textbf{恒转矩调速}；低频时定子电阻压降占比大，须做\textbf{低频电压补偿}（$U_1$ 略抬高）；
  \item \textbf{基频以上}：电压不能超过额定值，只能升频不升压，磁通随频率升高而减小，实现\textbf{恒功率调速}（弱磁调速）；
  \item 高性能调速还需\textbf{转差频率控制}、\textbf{矢量控制}（把定子电流分解为励磁分量与转矩分量分别控制）或\textbf{直接转矩控制}（DTC），以获得快速转矩响应。
\end{itemize}

% ============================================================
\section{不间断电源}

\textbf{不间断电源}（UPS）是当交流输入电源（习惯称为市电）发生中断或异常时，能持续向负载供电的装置，广泛用于计算机系统、通信机房、医院手术室、工业控制系统等不能断电的场合。

\begin{itemize}
  \item \textbf{组成}：整流器（把市电变成直流并向蓄电池充电）、蓄电池组（储能）、逆变器（把直流变成高质量交流）、静态开关（市电与逆变输出的切换）。
  \item \textbf{后备式 UPS}：市电正常时由市电直接供电（经简单滤波），断电时切换到逆变供电，切换时间约几毫秒；结构简单、成本低、效率高，但输出波形与电压质量一般。
  \item \textbf{在线式 UPS}：市电经“整流—逆变”始终由逆变器供电，输出波形纯净、无切换时间，是最可靠的方案，但效率较低、成本高。
  \item \textbf{在线互动式 UPS}：介于两者之间，市电正常时经变压器（可调压）供电，异常时逆变供电。
\end{itemize}

小容量 UPS 的主电路中，整流部分常用\textbf{二极管整流 + 直流斩波器}（兼作 PFC），以降低成本并提高输入功率因数。

% ============================================================
\section{开关电源}

各种电子设备需要多路不同电压的直流供电（如数字电路 5 V、3.3 V、1.8 V 等）。早期采用“工频变压器降压 + 整流滤波 + 线性调压”的\textbf{线性电源}：结构简单、纹波小，但\textbf{效率低（$30\%\sim50\%$）、体积大、重量重}。

采用\textbf{高频开关方式}进行电能变换的电源称为\textbf{开关电源}：市电整流滤波后经高频 DC-DC 变换器（几十到几百千赫兹）降压，再用高频变压器隔离、输出滤波。

\subsection{开关电源的结构}

典型结构：输入滤波与整流 → 功率因数校正（可选）→ 高频逆变（半桥/全桥/正激/反激）→ 高频变压器隔离 → 输出整流滤波 → 反馈控制。

与线性电源相比的优点：\textbf{效率高（可达 90\% 以上）、体积小、重量轻、输入电压范围宽}；缺点：\textbf{输出纹波与噪声较大、电磁干扰强、电路复杂}。

\subsection{开关电源的控制方式}

\begin{itemize}
  \item \textbf{PWM 方式}：开关频率固定，通过改变占空比稳压——最常用；
  \item \textbf{PFM 方式}：导通脉宽固定，通过改变开关频率稳压——轻载效率高，但滤波器设计困难；
  \item \textbf{混合方式}：重载用 PWM、轻载用 PFM（或跳周期），兼顾效率与纹波。
\end{itemize}

典型的开关电源控制系统为\textbf{电压闭环}（采样输出电压与基准比较，经 PI 调节器产生占空比指令）；要求高的场合再加\textbf{电流内环}（电流模式控制），以改善动态响应与均流。

\subsection{开关电源的应用}

广泛用于计算机、通信设备、家用电器、仪器仪表、LED 照明、电动汽车充电器等。它是“电力电子技术进入信息电子领域”的桥梁——所有电子装置的电源部分都是电力电子电路。

% ============================================================
\section{功率因数校正技术}

传统的二极管整流 + 电容滤波电路，输入电流是\textbf{窄脉冲}，功率因数只有 0.6 左右、谐波电流大（见 3.4、3.5 节）。为满足电网谐波标准（如 IEC 61000-3-2），须采用\textbf{功率因数校正（PFC）}技术。

\subsection{功率因数校正电路的基本原理}

目标：使输入电流\textbf{跟随输入电压的波形}（同频同相的正弦波），从而使 $\nu=I_1/I\to1$、$\cos\varphi_1\to1$，即 $\lambda\to1$。

\begin{itemize}
  \item \textbf{无源 PFC}：在整流桥后加电感（或“填谷”电路），把窄脉冲电流展宽，结构简单、成本低，可将 PF 提高到 0.8 左右；
  \item \textbf{有源 PFC（APFC）}：在整流桥与负载之间插入一个\textbf{Boost 变换器}，用控制环使输入电流跟踪输入电压的正弦波形：
  \begin{enumerate}
    \item 电压外环：采样输出直流电压，与基准比较后经误差放大器得到\textbf{电流指令的幅值}；
    \item 乘法器：把该幅值与整流后的输入电压波形相乘，得到\textbf{正弦电流指令}；
    \item 电流内环：控制开关使输入电流跟踪该指令，于是输入电流为正弦且与电压同相。
  \end{enumerate}
  \item 效果：PF 可达 0.99 以上、输入电流 THD 小于 5\%，且输出电压稳定（Boost 的升压特性使输出约为 400 V 直流）。
\end{itemize}

\keyformula{\lambda=\frac{I_1}{I}\cos\varphi_1\to1\qquad(\text{APFC 的目标})}

\subsection{单级功率因数校正技术}

上述基于 Boost 的有源 PFC 属“两级”结构（PFC 级 + DC-DC 级），成本较高。\textbf{单级 PFC} 把 PFC 功能与 DC-DC 变换合并到一级完成：在反激等变换器的输入端加少量元件（如填谷电容、耦合电感），使输入电流在半个工频周期内近似跟随输入电压，用一个开关同时完成 PFC 与稳压。

\begin{itemize}
  \item 优点：器件少、成本低，适合小功率（如 LED 驱动、适配器）；
  \item 缺点：储能电容电压随负载与输入电压变化较大、PF 与 THD 指标不如两级方案，输出纹波较大。
\end{itemize}

% ============================================================
\section{电力电子技术在电力系统中的应用}

\subsection{高压直流输电}

\textbf{高压直流输电（HVDC）}在长距离、大容量、跨区域输电以及海底电缆输电中有明显优势：没有交流线路的电容充电电流与稳定极限问题，线路损耗小、可控性强。

\begin{itemize}
  \item 送端\textbf{整流站}把交流变成直流，受端\textbf{逆变站}把直流变回交流，两端均为晶闸管换流阀（相控换流）。
  \item 传统的晶闸管换流阀需要电网提供换相电压（电网换相），故不能向无源网络供电；\textbf{轻型直流输电（柔性直流）}采用全控型器件 IGBT 构成的电压型换流器，可向弱电网甚至无源网络供电，是有源配电网、海上风电并网的重要技术。
  \item 换流站会产生大量谐波与无功消耗，须配\textbf{交流滤波器}与\textbf{无功补偿装置}。
\end{itemize}

\subsection{无功功率控制}

无功功率直接影响电网电压水平与线路损耗，对无功功率的控制是电力系统运行的核心问题之一。电力电子型无功补偿装置：

\begin{itemize}
  \item \textbf{晶闸管控制电抗器（TCR）}：用晶闸管的相控调节电抗器电流，连续吸收感性无功；
  \item \textbf{晶闸管投切电容器（TSC）}：分组投切电容器，分级提供容性无功（在电压过零投入以避免冲击）；
  \item \textbf{静止无功补偿器（SVC）}：TCR 与 TSC（或固定电容器 FC）组合，可连续调节；
  \item \textbf{静止无功发生器（SVG / STATCOM）}：用全控型器件构成的电压型逆变器，输出与电网同频的电压，通过调节其幅值控制无功电流的方向与大小，响应快、谐波小、低压时仍能输出额定无功，性能优于 SVC。
\end{itemize}

\subsection{谐波抑制与电能质量控制}

\begin{itemize}
  \item \textbf{无源滤波器}：LC 谐振支路，对特定次谐波呈低阻抗，结构简单但只能滤固定次谐波、易与电网谐振；
  \item \textbf{有源电力滤波器（APF）}：检测负载谐波电流，用 PWM 逆变器产生\textbf{大小相等、方向相反}的补偿电流将其抵消，可动态补偿各次谐波与无功，是治理谐波最有效的手段；
  \item \textbf{统一电能质量调节器、动态电压恢复器（DVR）}：用于抑制电压跌落、闪变、瞬时中断，改善供电质量；
  \item \textbf{柔性交流输电（FACTS）}：用电力电子装置控制线路的阻抗、电压与相位，提高输电能力与稳定性。
\end{itemize}

% ============================================================
\section{电力电子技术的其他应用}

\begin{itemize}
  \item \textbf{焊接电源}：逆变式焊机把工频整流后经高频逆变（20 kHz 以上）再整流输出，变压器体积重量大幅减小、动态响应快、焊接工艺性好，已取代传统工频焊机。
  \item \textbf{照明电源}：电子镇流器（高频点亮荧光灯，提高光效、消除频闪）、LED 驱动电源（恒流输出，通常含 PFC 级）；电力电子照明电源被称为“节能灯”，正逐步取代白炽灯与荧光灯。
  \item \textbf{电解与电镀}：需要低电压大电流的直流电源，多用双反星形或十二脉波整流（见 3.6 节）。
  \item \textbf{感应加热}：中频/高频电源（并联谐振式逆变电路，见 4.3 节）用于熔炼、淬火、焊接。
  \item \textbf{新能源与储能}：光伏并网逆变器、风力发电变流器、电池储能 PCS、电动汽车充电机——都是电力电子装置，且都要求高效率、高功率密度、单位功率因数与能量双向流动。
\end{itemize}

\corepoint{本章复习抓“装置—电路—指标”的对应关系：\textbf{直流调速}对应相控整流 + 电动机机械特性（连续/断续、整流/逆变）；\textbf{变频器}对应交直交（电压型/电流型、能否再生）与 $U_1/f_1=$ 常数；\textbf{UPS} 对应后备式/在线式与静态开关；\textbf{开关电源}对应高频 DC-DC + PWM 控制；\textbf{PFC} 对应“输入电流跟随输入电压”的 Boost APFC 与乘法器控制；\textbf{电力系统}对应 HVDC、SVC/SVG、APF。答题时先说\textbf{用什么电路}，再说\textbf{靠什么控制量调节}，最后说\textbf{优缺点与适用场合}。}
